\chapter{Literature Review}

\section{Classical economic dispatch}
The classical treatment of ELD is given by Wood, Wollenberg and Sheblé \cite{wood2014}. With quadratic costs $C_i(P_i) = a_iP_i^2 + b_iP_i + c_i$, the optimality condition is that every unit not at a limit operates at the same incremental cost $\lambda = dC_i/dP_i$; the $\lambda$-iteration method searches for this common value. Transmission losses are commonly represented by Kron's loss formula $\PL = P^\top B P + B_0^\top P + B_{00}$, whose coefficients are derived from the network. General nonlinear programming solvers, such as SLSQP \cite{kraft1988}, solve the same problem directly and handle additional constraints.

\section{Heuristic methods for ELD}
Population-based heuristics are widely used for non-convex dispatch. Gaing \cite{gaing2003} applied particle swarm optimisation (PSO) to ELD with generator constraints and handled the power-balance equality by computing one ``dependent'' unit from the others, which keeps every candidate exactly balanced. The inertia-weight PSO of Shi and Eberhart \cite{shi1998} is the usual variant. Real-coded genetic algorithms with simulated binary crossover (SBX) \cite{deb1995} and simulated annealing \cite{kirkpatrick1983} are other standard choices. These methods are used in this project as benchmarks.

\section{QUBO and Ising formulations}
Lucas \cite{lucas2014} showed how many NP problems can be written as Ising models: binary decision variables become spins, objectives become linear and quadratic terms, and constraints become quadratic penalties. The substitution $x = (1 - z)/2$ converts a QUBO in binary variables $x \in \{0,1\}$ into an Ising Hamiltonian in spins $z \in \{-1,+1\}$. This is the formulation route followed in this project.

\section{QAOA}
QAOA \cite{farhi2014} alternates a cost unitary $e^{-i\gamma H_C}$ and a mixer unitary $e^{-i\beta H_M}$ for $p$ layers, starting from the uniform superposition. Zhou et al. \cite{zhou2020} studied QAOA performance and proposed the INTERP strategy, in which optimised angles at depth $p$ are interpolated to initialise depth $p + 1$, avoiding poor random starts. Barkoutsos et al. \cite{barkoutsos2020} proposed training with the conditional value-at-risk (CVaR) of the sampled energies instead of the mean, which focuses the optimisation on the best outcomes and improves the probability of sampling the optimum.

\section{Quantum methods for power-system dispatch}
Han, Li and Xu \cite{han2024} proposed a quantum-assisted method for stochastic economic dispatch in renewable-rich systems. Jing, Wang and Li \cite{jing2023} developed a data-driven variant of QAOA for power-system optimisation problems. Wang, Cao and Li \cite{wang2024} applied a quantum-behaved particle swarm optimiser, a quantum-inspired classical heuristic, to economic dispatch. These works motivate the use of quantum algorithms in power systems but generally do not report a full comparison against exact and heuristic baselines on a common criterion, or AC validation of the resulting dispatch.

\section{Research gap}
From the review, the following gaps are addressed in this project:
\begin{itemize}
\item A transparent ELD-to-QUBO formulation with a justified penalty weight rather than an arbitrary one.
\item Measurement of QAOA accuracy against the exact optimum over a full daily load profile.
\item A fair comparison with standard algorithms using the same success criterion and evaluation budget.
\item Validation of the quantum dispatch with an AC load flow, including losses, voltages and line loading.
\item An honest report of real-hardware performance and of the absence of a speed-up.
\end{itemize}
